% Source: source/普通高考/2013/2013广东理.pdf, page 1, question 11.
% pdfimages -list: no embedded image objects; the program flowchart is vector line art.
% Grid evidence: tmp/review_2013_guangdong_science/grid/page_grid100.png,
%   tmp/review_2013_guangdong_science/grid/page_grid50.png.
% No-grid source crop: tmp/review_2013_guangdong_science/crop/q11_orig_fig_final.png.
% Node→directed-edge list (verified against the source page):
% 开始→输入 n→i=1,s=1→i≤n?; 是→s=s+(i−1)→i=i+1→(loop back to i≤n?);
% 否→输出 s→结束.  The loop branch is independent of the output branch.
% solid segments: all node borders and connector shafts; dashed segments: none.
% labels: formulas and Chinese text are centered in each node with local zero indentation;
% 是 is beside the true downward edge and 否 beside the false right edge.

\begin{tikzpicture}[
  x=1cm,
  y=0.92cm,
  >=Stealth,
  line width=0.45pt,
  line cap=round,
  line join=round,
  font=\normalsize,
  every node/.style={anchor=center,align=center,inner sep=0pt,
    execute at begin node={\setlength{\parindent}{0pt}}},
  flow/.style={draw,anchor=center,align=center,inner sep=0pt,
    execute at begin node={\setlength{\parindent}{0pt}}},
  terminal/.style={flow,rounded corners=0.16cm,minimum width=1.10cm,minimum height=0.64cm},
  process/.style={flow,minimum width=2.50cm,minimum height=0.68cm,inner xsep=4pt,inner ysep=2pt},
  io/.style={flow,shape=trapezium,trapezium left angle=72,trapezium right angle=108,
    minimum width=1.90cm,minimum height=0.76cm,inner xsep=4pt,inner ysep=2pt},
  decision/.style={flow,diamond,aspect=2.75,minimum width=2.75cm,minimum height=0.86cm,
    inner xsep=4pt,inner ysep=2pt},
  branchlabel/.style={anchor=center,align=center,inner sep=0pt,
    execute at begin node={\setlength{\parindent}{0pt}}},
]
  \path[use as bounding box] (-2.15,1.75) rectangle (5.35,9.62);

  \node[terminal] (start) at (0,9.12) {开始};
  \node[io] (input) at (0,7.94) {输入 $n$};
  \node[process,minimum width=2.72cm] (init) at (0,6.70) {$i=1,s=1$};
  \node[decision,minimum width=3.05cm,minimum height=0.98cm] (test) at (0,5.32) {$i\leq n$};
  \node[process,minimum width=2.95cm] (sum) at (0,3.90) {$s=s+(i-1)$};
  \node[process,minimum width=2.38cm] (inc) at (0,2.62) {$i=i+1$};
  \node[io,minimum width=1.72cm] (output) at (4.20,4.48) {输出 $s$};
  \node[terminal] (finish) at (4.20,3.20) {结束};

  \draw[->] (start.south) -- (input.north);
  \draw[->] (input.south) -- (init.north);
  \draw[->] (init.south) -- (test.north);
  \draw[->] (test.south) -- node[branchlabel,left=2pt,pos=0.18] {是} (sum.north);
  \draw[->] (sum.south) -- (inc.north);

  % False branch: right exit to output and then termination.
  \draw[->] (test.east) -- node[branchlabel,above=2pt] {否} (4.20,5.32) -- (output.north);
  \draw[->] (output.south) -- (finish.north);

  % Loop-back branch: i=i+1 returns to the decision without touching output.
  \coordinate (loop-bottom) at (0,2.05);
  \coordinate (loop-left) at (-2.00,2.05);
  % Keep the return leg at least 2.5 mm clear of the initialization box.
  \coordinate (loop-top) at (-2.00,6.00);
  \coordinate (loop-junction) at (0,6.00);
  % The source return exits below i=i+1, runs left, then rises to the main stem.
  \draw (inc.south) -- (loop-bottom) -- (loop-left) -- (loop-top);
  % The arrowhead stops before the shared stem; the remaining 4 mm tail is
  % unarrowed, keeping it clear of the following downward arrowhead.
  \draw[->] (loop-top) -- (-0.40,6.00);
  \draw (-0.40,6.00) -- (loop-junction);
\end{tikzpicture}
